\documentclass[11pt,a4paper]{article}

% Template ACL chinh thuc
\usepackage[final]{acl}

% Ho tro tieng Viet voi XeLaTeX
\usepackage{fontspec}
\setmainfont{Noto Serif}
\usepackage{microtype}

\title{
Xây dựng và so sánh hai mô hình dịch máy Anh--Việt:
Seq2Seq LSTM với Attention và Transformer
}

% THAY HO TEN VA MSSV TRUOC KHI NOP
\author{
Nguyễn Tất Tú: 24022481 \quad
Nguyễn Văn Trung: 24022475 \quad
Dương Hoàng Đức: 24022290 \quad \\
Nhóm: 1
}

\begin{document}
\maketitle

\begin{abstract}
Đề tài tự cài đặt bằng PyTorch và so sánh hai mô hình
dịch máy Anh--Việt, Seq2Seq LSTM với Attention và
Transformer, trên IWSLT'15 với cùng tokenizer, giải mã
và cách đánh giá. Nhóm so sánh SacreBLEU, chi phí huấn
luyện và hiệu quả theo độ dài câu, đồng thời trực quan
hóa attention và phân tích lỗi.
\end{abstract}

\section{Giới thiệu}

Dịch máy thần kinh (Neural Machine Translation,
NMT) chuyển một chuỗi ngôn ngữ nguồn thành chuỗi
ngôn ngữ đích bằng mạng học sâu. Kiến trúc Seq2Seq
dùng encoder và decoder; khi kết hợp với Attention,
decoder tập trung vào những vị trí liên quan trong
câu nguồn \citep{bahdanau2015nmt}. Transformer thay
cơ chế hồi quy bằng self-attention, giúp huấn luyện
song song hiệu quả hơn \citep{vaswani2017attention}.

Mục tiêu của đề tài là xây dựng và so sánh hai mô
hình trên cùng bài toán dịch Anh sang Việt, tập trung
vào chất lượng dịch, chi phí tính toán và các lỗi
thường gặp. Chúng tôi kỳ vọng Transformer cho BLEU
cao hơn và huấn luyện nhanh hơn mỗi epoch, trong khi
LSTM có thể cạnh tranh trên câu ngắn.

\section{Dữ liệu}

Nghiên cứu dùng IWSLT'15 English--Vietnamese
\citep{luong2015iwslt} gồm 133.317 cặp câu huấn luyện,
1.553 cặp validation (\texttt{tst2012}) và 1.268 cặp
kiểm thử (\texttt{tst2013}). Nhóm giữ nguyên cách chia.
Validation dùng để chọn siêu tham số và checkpoint;
test chỉ dùng cho đánh giá cuối.

Tiền xử lý gồm chuẩn hóa Unicode NFC (giữ nguyên chữ
hoa), loại cặp câu dài quá 100 token hoặc có tỉ lệ độ
dài nguồn/đích lệch quá 2,5 lần, thêm token đặc biệt
và padding. Tokenizer mặc định là SentencePiece BPE
huấn luyện chung trên phần train của hai ngôn ngữ, từ
vựng 8.000 token, dùng chung embedding giữa encoder
và decoder.

\section{Phương pháp}

\paragraph{Phạm vi.} \emph{Bắt buộc:} hai mô hình,
SacreBLEU, attention, ablation (Mục~\ref{sec:exp}).
\emph{Mở rộng, theo thứ tự nếu còn thời gian:}
(1) so sánh tokenizer BPE, SentencePiece, mức từ;
(2) fine-tune EnViT5 hoặc mBART làm mốc tham chiếu.

\paragraph{Mô hình 1: Seq2Seq LSTM với Attention.}
Encoder BiLSTM biểu diễn câu tiếng Anh, decoder LSTM
sinh lần lượt các token tiếng Việt. Tại mỗi bước, Attention
kiểu Luong (general) tính trọng số giữa trạng thái decoder
và các trạng thái encoder để tạo context vector, kết hợp
input feeding.

\paragraph{Mô hình 2: Transformer.}
Transformer Encoder--Decoder nhỏ theo
\citet{vaswani2017attention}, tự cài đặt embedding,
positional encoding, multi-head attention, feed-forward
network, residual connection, layer normalization;
decoder dùng causal mask. Mô hình được thu nhỏ so với
bản gốc để tránh overfit trên dữ liệu 133 nghìn cặp câu.

\begin{table*}[t]
\centering
\footnotesize
\begin{tabular}{lll}
\hline
& \textbf{Seq2Seq LSTM + Attention} & \textbf{Transformer nhỏ} \\
\hline
Encoder & BiLSTM 2 lớp, hidden 512 mỗi chiều & 3 lớp \\
Decoder & LSTM 2 lớp, hidden 512 & 3 lớp, causal mask \\
Attention & Luong (general) & Multi-head, 4 head \\
Kích thước & embedding 512 & $d_{model}=512$, $d_{ff}=1024$ \\
Dropout & 0,3 & 0,3 \\
Optimizer & Adam, lr $10^{-3}$, gradient clip 1,0 &
Adam ($\beta_2=0{,}98$), lr đỉnh $5\cdot10^{-4}$, warmup 2.000 bước \\
Chung & \multicolumn{2}{l}{label smoothing 0,1; embedding decoder dùng chung với lớp đầu ra; mixed precision} \\
Tham số (ước tính) & $\sim$22M & $\sim$20M \\
\hline
\end{tabular}
\caption{Cấu hình dự kiến. Số tham số sẽ được đo lại bằng mã nguồn.}
\label{tab:config}
\end{table*}

\paragraph{Tiêu chí công bằng.} Hai mô hình cùng dữ liệu,
tokenizer, giải mã, GPU và ngân sách tối đa 30 epoch với
early stopping (kiên nhẫn 5). So sánh chính theo cùng cỡ
tham số (Bảng~\ref{tab:config}); thời gian huấn luyện chỉ
được báo cáo kèm.

\section{Thiết kế thực nghiệm}
\label{sec:exp}

Hai mô hình huấn luyện bằng PyTorch trên GPU miễn phí
(Kaggle hoặc Colab); nhóm theo dõi training/validation
loss và lưu checkpoint tốt nhất trên validation.

\paragraph{Đánh giá.} Corpus BLEU \citep{papineni2002bleu}
qua SacreBLEU (tokenization mặc định) trên văn bản đã
detokenize, tập \texttt{tst2013} (không so sánh trực tiếp
với bài báo dùng quy ước khác). Giải mã bằng beam search
(beam 5, chuẩn hóa độ dài, độ dài tối đa bằng độ dài
nguồn cộng 50), dùng chung cho cả hai mô hình. Mỗi mô
hình chạy một seed; nếu còn thời gian, ước lượng khoảng
tin cậy bằng paired bootstrap.

\paragraph{So sánh.} BLEU tổng thể; BLEU theo bốn nhóm độ
dài nguồn ($\le$10, 11--20, 21--30, $>$30 từ); số tham số,
thời gian mỗi epoch, tổng thời gian đến checkpoint tốt
nhất và bộ nhớ GPU đỉnh.

\paragraph{Ablation.} (A1) greedy so với beam search;
(A2) Transformer 2 so với 3 lớp; (A3) LSTM có và không
có Attention.

\paragraph{Attention và phân tích lỗi.} Với LSTM, vẽ trọng
số attention của các bước giải mã. Với Transformer, vẽ
cross-attention lớp decoder cuối: trung bình các head và
từng head cho 2--3 câu minh họa. Phân tích lỗi trên 50 câu
test (phân tầng theo độ dài), gán vào bốn nhóm: thiếu hoặc
lặp từ, sai từ vựng, sai trật tự, sai tên riêng hoặc số.

\section{Phân công và rủi ro}

Tiến độ theo timeline chung của môn học; dưới đây là
phân công và phương án dự phòng:

\noindent\textbf{Nguyễn Tất Tú:} dữ liệu IWSLT'15, tokenizer,
DataLoader, BLEU và Proposal.

\noindent\textbf{Nguyễn Văn Trung:} Seq2Seq LSTM Attention, beam search dùng chung,
demo và báo cáo tiến độ.

\noindent\textbf{Dương Hoàng Đức (nhóm trưởng):} Transformer, huấn luyện,
checkpoint, logging và báo cáo cuối kỳ.

\textbf{Rủi ro.}
(i) Giới hạn phiên Colab/Kaggle: lưu checkpoint mỗi epoch. (ii) Transformer overfit hoặc hội tụ chậm: tăng dropout, giảm xuống 2 lớp hoặc tăng warmup. (iii) Chậm tiến độ: bỏ EnViT5/mBART trước, rồi so sánh tokenizer; giữ nguyên phần bắt buộc.

\begingroup
\small
\setlength{\bibsep}{1pt}
\bibliographystyle{acl_natbib}
\bibliography{references}
\endgroup

\end{document}